% ============================================================
%  SEZIONE 5 — Contributo dei componenti del gruppo
% ============================================================
\section{Contributo di ogni componente del gruppo}
% La Tabella~\ref{tab:contributi} riassume le attività svolte da ciascun
% componente del team.

\begin{table}[H]
\centering
% \caption{Ripartizione del lavoro tra i membri del team.}
\label{tab:contributi}
\renewcommand{\arraystretch}{1.3}
\begin{tabularx}{\linewidth}{l X}
\toprule
\textbf{Componente} & \textbf{Attività svolte} \\
\midrule
\textbf{Stefano Videsott} &
  Architettura iniziale della web app principale; Spring Security,
  controller, service layer e repository; parte del frontend frontend (relativo alle parti sviluppate), autenticazione e registrazione, dahsboards con alcune delle funzionalità, logica legata alle recensioni con
  \texttt{fetch}/\texttt{async}/\texttt{await}; statistiche Chart.js;
  refactoring completo finale della codebase per
  uniformare i contributi di entrambi i membri (nomi di classi,
  metodi, variabili, tabelle e colonne DB rinominati in inglese con
  convenzioni coerenti). \\
\textbf{Alessandro Como} &
  Realizzazione della homepage e della pagina di contatto. Progettazione e sviluppo del REST service: controller, repository,
  schema SQL e dati di test; integrazione con OpenFeign REST;
  buisness logic e view relative alla visualizzazione e al completamento degli allenamenti e delle parti collegate. \\
\bottomrule
\end{tabularx}
\end{table}

La prima parte di ideazione e modellazione dello schema del database è stata fatta insieme.
